\paragrafos{Βασικές έννοιες: σημείο, ευθεία, ημιευθεία, ευθύγραμμο τμήμα, γωνία, κύκλος}{Β.1.1, «Βασικές έννοιες: σημείο, ευθεία, ημιευθεία, ευθύγραμμο τμήμα, γωνία, κύκλος»}

\begin{anagke}
Η κυρία Αλεξάνδρα ανακοινώνει ένα μεγάλο εγχείρημα για τη Γεωμετρία της χρονιάς: «Ο κήπος του σχολείου θα ανανεωθεί και η Α3 θα τον σχεδιάσει. Θα φέρουμε τα δέντρα, τα μονοπάτια, τον φράχτη, ένα στρογγυλό παρτέρι. Αλλά πρώτα πρέπει να μάθουμε να μιλάμε για αυτά που θα σχεδιάσουμε». Η Αγγελική, που βλέπει σχήματα παντού, χαίρεται: «Επιτέλους ένα μάθημα για αυτό που κάνω όλη μέρα!».

Η Αγγελική παίρνει ένα χαρτί και κάνει μια κάτοψη του κήπου. Κάθε δέντρο το σημειώνει με μια τελεία. Το μονοπάτι ανάμεσα σε δύο δέντρα το ενώνει με ένα τεντωμένο νήμα. Ο παππούς Γιάννης, που έρχεται να βοηθήσει, λέει: «Για να φυτέψετε σε ευθεία γραμμή, τεντώστε σχοινί και προχωρήστε όσο θέλετε προς τις δύο μεριές». Ο Στέλιος ρωτάει: «Μέχρι πού; Ο φράχτης τελειώνει στον τοίχο!». Η Ρήνη δείχνει τον φακό της: η δέσμη του φωτός ξεκινά από ένα σημείο και συνεχίζει προς μία μεριά χωρίς να τελειώνει. Για το στρογγυλό παρτέρι ο παππούς καρφώνει ένα παλούκι, δένει ένα σχοινί και γυρίζει τεντωμένο γύρω του. «Όλα τα σημεία της γραμμής απέχουν το ίδιο από το παλούκι», παρατηρεί η Φανή και το γράφει στο τετράδιό της.

Κάθε μία από αυτές τις εικόνες είναι μια γεωμετρική έννοια: σημείο, ευθύγραμμο τμήμα, ευθεία, ημιευθεία, γωνία, κύκλος. Σε αυτή την παράγραφο τις ονομάζουμε προσεκτικά και μαθαίνουμε να τις γράφουμε με σύμβολα.
\end{anagke}

\begin{lexi}{σημείο, ευθεία, επίπεδο, κύκλος}{σημείο, ευθεία, επίπεδο, κύκλος}
Το \lexiword{σημείο} είναι «αυτό που σημαδεύεται»: ένα σημάδι χωρίς διαστάσεις. Η \lexiword{ευθεία} πηγαίνει «ευθύ», χωρίς να στρίβει. Το \lexiword{επίπεδο} είναι ό,τι είναι «επί του πέδου», ίσιο σαν το δάπεδο. Το \lexiword{ευθύγραμμο τμήμα} είναι ένα «κομμένο» κομμάτι ευθείας. Το \lexiword{ημι-} σημαίνει «μισό»: ημιευθεία, ημιεπίπεδο. Ο \lexiword{κύκλος} είναι ό,τι «κυλάει» και γυρίζει. Η \lexiword{ακτίνα} είναι η «αχτίδα» που ξεκινά από το κέντρο.
\end{lexi}

\begin{metafrasi}
Το ίχνος της μύτης του μολυβιού. & σημείο $\mathrm{Α}$\\
Τεντωμένη κλωστή με άκρα $\mathrm{Α}$ και $\mathrm{Β}$. & ευθύγραμμο τμήμα $\mathrm{ΑΒ}$\\
Το ίδιο τμήμα προεκτεινόμενο απεριόριστα και προς τις δύο μεριές. & ευθεία $\mathrm{ΑΒ}$ ή $\varepsilon$ ή $x'x$\\
Φως φακού που ξεκινά από ένα σημείο. & ημιευθεία $\mathrm{Αx}$\\
Δύο ημιευθείες με κοινή αρχή $\mathrm{Ο}$. & γωνία $x\hat{\mathrm{Ο}}y$\\
Τα σημεία που απέχουν $\rho$ από το $\mathrm{Ο}$. & κύκλος $(\mathrm{Ο},\rho)$\\
\end{metafrasi}

\tmima{Η θεωρία}

\tmimasmall{Σημείο και ευθεία}

Ένα \textbf{σημείο} παριστάνεται με μία τελεία και δεν έχει διαστάσεις (ούτε μήκος ούτε πλάτος). Το ίχνος της μύτης του μολυβιού μας δίνει την έννοια του σημείου. Το συμβολίζουμε με ένα κεφαλαίο γράμμα: $\mathrm{Α}$, $\mathrm{Β}$, $\Gamma$.

\textbf{Ευθεία} είναι η γραμμή που εκτείνεται απεριόριστα, χωρίς αρχή και τέλος, και είναι απολύτως ίσια. Τη συμβολίζουμε με ένα μικρό γράμμα, π.χ. $\varepsilon$, ή με δύο γράμματα, π.χ. $x'x$.

\begin{idiotita}{ευθεία και σημεία}
Από ένα σημείο διέρχονται \textbf{άπειρες} ευθείες. Από δύο διαφορετικά σημεία διέρχεται \textbf{μία και μοναδική} ευθεία.
\end{idiotita}

\begin{center}
\begin{tikzpicture}[scale=0.9]
\begin{scope}
\foreach \a in {0,35,70,105,140,-35} {\draw[gblue] ({-1.5*cos(\a)},{-1.5*sin(\a)})--({1.5*cos(\a)},{1.5*sin(\a)});}
\fill (0,0) circle (2pt) node[below right] {$\mathrm{Α}$};
\node at (0,-1.9) {\small άπειρες ευθείες};
\end{scope}
\begin{scope}[xshift=5.5cm]
\draw[gblue,{Stealth}-{Stealth}] (-1.7,0)--(1.7,0);
\fill (-0.8,0) circle (2pt) node[above] {$\mathrm{Α}$}; \fill (0.8,0) circle (2pt) node[above] {$\mathrm{Β}$};
\node at (0,-0.7) {\small μία ευθεία};
\end{scope}
\end{tikzpicture}
\end{center}

\tmimasmall{Ευθύγραμμο τμήμα και ημιευθεία}

Αν σε μια ευθεία πάρουμε δύο σημεία $\mathrm{Α}$ και $\mathrm{Β}$, το σχήμα που αποτελείται από τα $\mathrm{Α}$, $\mathrm{Β}$ και όλα τα σημεία ανάμεσά τους λέγεται \textbf{ευθύγραμμο τμήμα} $\mathrm{ΑΒ}$ (ή $\mathrm{ΒΑ}$). Τα $\mathrm{Α}$ και $\mathrm{Β}$ είναι τα \textbf{άκρα} του. Το τμήμα έχει αρχή και τέλος, ενώ η ευθεία όχι.

Αν προεκτείνουμε το τμήμα $\mathrm{ΑΒ}$ και προς τα δύο άκρα του, σχηματίζεται ευθεία. Αν το προεκτείνουμε \emph{μόνο προς το ένα} άκρο, το σχήμα έχει αρχή αλλά δεν έχει τέλος και λέγεται \textbf{ημιευθεία}. Τη συμβολίζουμε $\mathrm{Αx}$, με πρώτο το γράμμα της αρχής.

\begin{center}
\begin{tikzpicture}[scale=0.95]
\draw[thick] (0,0)--(3,0); \fill (0,0) circle (2pt) node[below] {$\mathrm{Α}$}; \fill (3,0) circle (2pt) node[below] {$\mathrm{Β}$};
\node at (1.5,-0.9) {τμήμα $\mathrm{ΑΒ}$};
\begin{scope}[xshift=4.8cm]
\draw[thick,{Stealth}-{Stealth}] (-0.5,0)--(3.5,0); \fill (0.5,0) circle (2pt) node[below] {$\mathrm{Α}$}; \fill (2.7,0) circle (2pt) node[below] {$\mathrm{Β}$}; \node[above] at (3.2,0) {$\varepsilon$};
\node at (1.5,-0.9) {ευθεία $\varepsilon$};
\end{scope}
\begin{scope}[xshift=9.8cm]
\draw[thick,-{Stealth}] (0,0)--(3.2,0) node[above] {$x$}; \fill (0,0) circle (2pt) node[below] {$\mathrm{Α}$};
\node at (1.5,-0.9) {ημιευθεία $\mathrm{Αx}$};
\end{scope}
\end{tikzpicture}
\end{center}

Ένα σημείο $\mathrm{Α}$ μιας ευθείας $x'x$ τη χωρίζει σε δύο ημιευθείες, τις $\mathrm{Αx}$ και $\mathrm{Αx'}$, που έχουν κοινή αρχή το $\mathrm{Α}$. Οι ημιευθείες αυτές λέγονται \textbf{αντικείμενες}.

\begin{center}
\begin{tikzpicture}
\draw[thick,{Stealth}-{Stealth}] (-3,0)--(3,0); \fill (0,0) circle (2pt) node[above] {$\mathrm{Α}$};
\node[above] at (-2.8,0) {$x'$}; \node[above] at (2.8,0) {$x$};
\node[below,gblue] at (-1.5,0) {\small $\mathrm{Αx'}$}; \node[below,gblue] at (1.5,0) {\small $\mathrm{Αx}$};
\end{tikzpicture}
\end{center}

\tmimasmall{Γωνία}

Σχεδιάζουμε δύο ημιευθείες $\mathrm{Οx}$ και $\mathrm{Οy}$ με κοινή αρχή το σημείο $\mathrm{Ο}$. Οι ημιευθείες χωρίζουν το επίπεδο σε δύο περιοχές. Καθεμία από αυτές τις περιοχές, μαζί με τις δύο ημιευθείες, ονομάζεται \textbf{γωνία}. Το $\mathrm{Ο}$ είναι η \textbf{κορυφή} και οι ημιευθείες οι \textbf{πλευρές} της γωνίας.
\begin{itemize}[leftmargin=1.2em,itemsep=0.3ex,topsep=2pt]
\item Η μικρότερη περιοχή δίνει την \textbf{κυρτή} γωνία.
\item Η μεγαλύτερη περιοχή δίνει τη \textbf{μη κυρτή} γωνία.
\end{itemize}

\begin{center}
\begin{tikzpicture}[scale=1.1]
\fill[gblueL] (0,0)--(2.5,0) arc(0:38:2.5)--cycle;
\draw[thick] (2.9,0)--(0,0)--(2.3,1.8); \fill (0,0) circle (2pt) node[below left] {$\mathrm{Ο}$};
\node[right] at (2.9,0) {$x$}; \node[above right] at (2.3,1.8) {$y$};
\draw[gblue] (0.8,0) arc(0:38:0.8); \node[gblue] at (1.15,0.33) {\small κυρτή};
\draw[ggray] (-0.6,0) arc(180:322:0.6); \node[ggray] at (-1.15,-0.6) {\small μη κυρτή};
\end{tikzpicture}
\end{center}

Συμβολίζουμε τη γωνία με τρία γράμματα, γράφοντας στη μέση το γράμμα της κορυφής: $x\hat{\mathrm{Ο}}y$ ή $y\hat{\mathrm{Ο}}x$. Μπορούμε επίσης να τη συμβολίσουμε $\hat{\mathrm{Ο}}$ (με το γράμμα της κορυφής) ή με ένα μικρό ελληνικό γράμμα, π.χ. $\hat{\omega}$.

\textbf{Σημείωση.} Με τον γνώμονα κατασκευάζουμε ευθείες που λέγονται \textbf{κάθετες}. Η γωνία που σχηματίζεται λέγεται \textbf{ορθή}.

\tmimasmall{Επίπεδο και ημιεπίπεδο}

\textbf{Επίπεδο} είναι η επιφάνεια πάνω στην οποία εφαρμόζει παντού η ευθεία γραμμή. Το επίπεδο επεκτείνεται απεριόριστα και το ονομάζουμε με ένα κεφαλαίο γράμμα, π.χ. $\Pi$. Κάθε ευθεία $\varepsilon$ ενός επιπέδου το χωρίζει σε δύο \textbf{ημιεπίπεδα}.

\begin{center}
\begin{tikzpicture}[scale=0.95]
\fill[gblueL] (0,0)--(4.4,0)--(5.4,1.8)--(1,1.8)--cycle;
\fill[gwarm] (0,0)--(4.4,0)--(3.9,-0.9)--(-0.5,-0.9)--cycle;
\draw (-0.5,-0.9)--(3.9,-0.9)--(5.4,1.8)--(1,1.8)--cycle;
\draw[thick,gblue] (-0.25,-0.45)--(5.15,1.35) coordinate (e);
\node[gblue] at (5.5,1.5) {$\varepsilon$};
\node at (1.4,1.1) {$\Pi_1$}; \node at (3.5,-0.45) {$\Pi_2$};
\end{tikzpicture}
\end{center}

\tmimasmall{Κύκλος}

\begin{orismos}{κύκλος}
\textbf{Κύκλος} λέγεται το σύνολο όλων των σημείων του επιπέδου που απέχουν την ίδια απόσταση από ένα σταθερό σημείο $\mathrm{Ο}$.
\end{orismos}

Το σημείο $\mathrm{Ο}$ λέγεται \textbf{κέντρο} του κύκλου. Η σταθερή απόσταση λέγεται \textbf{ακτίνα} του κύκλου και συμβολίζεται με $\rho$. Έναν κύκλο με κέντρο $\mathrm{Ο}$ και ακτίνα $\rho$ τον γράφουμε $(\mathrm{Ο},\rho)$. Τον σχεδιάζουμε με τον \textbf{διαβήτη}. Οι κύκλοι που έχουν το ίδιο κέντρο λέγονται \textbf{ομόκεντροι}.

\begin{center}
\begin{tikzpicture}
\draw[thick] (0,0) circle (1.2); \draw[thick] (0,0) circle (1.9);
\fill (0,0) circle (2pt) node[below left] {$\mathrm{Ο}$};
\draw[gblue,thick] (0,0)--(0.85,0.85) node[midway,above left] {$\rho$};
\fill (0.85,0.85) circle (2pt) node[above right] {$\mathrm{Α}$};
\node at (0,-2.4) {\small δύο ομόκεντροι κύκλοι};
\end{tikzpicture}
\end{center}

\begin{prosoxi}
\begin{itemize}[leftmargin=1.2em,itemsep=0.3ex,topsep=2pt]
\item Το τμήμα $\mathrm{ΑΒ}$ και το $\mathrm{ΒΑ}$ είναι το ίδιο τμήμα. Αλλά οι ημιευθείες $\mathrm{Αx}$ και $\mathrm{Βx}$ είναι διαφορετικές, γιατί έχουν διαφορετική αρχή.
\item Το τμήμα έχει άκρα και μήκος. Η ευθεία και η ημιευθεία δεν έχουν δύο άκρα και δεν έχουν μήκος.
\item Στο όνομα της γωνίας το γράμμα της κορυφής γράφεται πάντα στη μέση: $x\hat{\mathrm{Ο}}y$, όχι $\mathrm{Ο}xy$.
\item Ο κύκλος είναι μόνο η γραμμή, όχι το εσωτερικό του. Όλες οι ακτίνες του ίδιου κύκλου είναι ίσες.
\end{itemize}
\end{prosoxi}

\tmima{Λυμένα παραδείγματα}

\paradeigma{Πόσες ευθείες;}
\tixero Τέσσερα δέντρα $\mathrm{Α},\,\mathrm{Β},\,\Gamma,\,\Delta$ του κήπου, από τα οποία τρία οποιαδήποτε δεν βρίσκονται στην ίδια ευθεία.\par
\tithelo Πόσες ευθείες ορίζουν.\par
\possindeo Κάθε ζεύγος σημείων ορίζει μία μοναδική ευθεία. Τα ζεύγη είναι $\mathrm{ΑΒ},\,\mathrm{ΑΓ},\,\mathrm{ΑΔ},\,\mathrm{ΒΓ},\,\mathrm{ΒΔ},\,\Gamma\Delta$. Με τον τρόπο των χειραψιών: $\dfrac{4\cdot3}{2}=6$.\par
\apantisi $6$ ευθείες.\par
\elegxos Αν το $\mathrm{Δ}$ ήταν πάνω στην ευθεία $\mathrm{ΑΒ}$, οι ευθείες θα ήταν λιγότερες.

\paradeigma{Ονόματα ημιευθειών}
\tixero Το σημείο $\mathrm{Ο}$ βρίσκεται πάνω στην ευθεία $x'x$. Τα $\mathrm{Α}$ και $\mathrm{Β}$ είναι δεξιά και αριστερά του $\mathrm{Ο}$ αντίστοιχα.\par
\tithelo Ποιες είναι οι δύο ημιευθείες που ορίζει το $\mathrm{Ο}$ και ποια περιέχει το $\mathrm{Α}$.\par
\possindeo Το $\mathrm{Ο}$ χωρίζει την ευθεία σε δύο ημιευθείες με αρχή το $\mathrm{Ο}$. Η μία πηγαίνει προς το $x$ και η άλλη προς το $x'$.\par
\apantisi Οι $\mathrm{Οx}$ (περιέχει το $\mathrm{Α}$) και $\mathrm{Οx'}$ (περιέχει το $\mathrm{Β}$), αντικείμενες.\par
\elegxos Αν τις ενώσω παίρνω ξανά ολόκληρη την ευθεία.

\paradeigma{Το όνομα μιας γωνίας}
\tixero Οι ημιευθείες $\mathrm{Οx}$ και $\mathrm{Οy}$ με κοινή αρχή $\mathrm{Ο}$ σχηματίζουν τη μικρότερη από τις δύο περιοχές του επιπέδου.\par
\tithelo Να ονομάσω τη γωνία με δύο τρόπους και να πω αν είναι κυρτή.\par
\possindeo Η κορυφή είναι το $\mathrm{Ο}$ και μπαίνει στη μέση: $x\hat{\mathrm{Ο}}y$ ή $y\hat{\mathrm{Ο}}x$. Επίσης μπορώ να γράψω απλώς $\hat{\mathrm{Ο}}$.\par
\apantisi Η γωνία $x\hat{\mathrm{Ο}}y$ είναι κυρτή, γιατί αντιστοιχεί στη μικρότερη περιοχή.\par
\elegxos Η μεγαλύτερη περιοχή θα έδινε τη μη κυρτή γωνία.

\paradeigma{Το μονοπάτι χωρίζει τον κήπο}
\tixero Ένα ευθύγραμμο μονοπάτι $\varepsilon$ χωρίζει τον κήπο σε δύο ημιεπίπεδα. Τα δέντρα $\mathrm{Κ}$ και $\Lambda$ είναι στο ένα και το $\mathrm{Μ}$ στο άλλο.\par
\tithelo Ποιες ευθύγραμμες διαδρομές ανάμεσα σε δύο δέντρα διασχίζουν το μονοπάτι.\par
\possindeo Για να πάω από το ένα ημιεπίπεδο στο άλλο πρέπει να περάσω από την ευθεία $\varepsilon$.\par
\apantisi Τα τμήματα $\mathrm{ΚΜ}$ και $\Lambda\mathrm{Μ}$ διασχίζουν το μονοπάτι, το $\mathrm{ΚΛ}$ όχι.\par
\elegxos Τα $\mathrm{Κ}$ και $\Lambda$ είναι στην ίδια πλευρά.

\paradeigma{Το στρογγυλό παρτέρι}
\tixero Το σχοινί του παππού είναι $1{,}5$ μ. και ο παππούς το κρατά τεντωμένο γύρω από το παλούκι $\mathrm{Ο}$. Ένα δέντρο $\mathrm{Δ}$ βρίσκεται $2$ μ. από το παλούκι.\par
\tithelo Ποιος κύκλος χαράσσεται και αν το δέντρο είναι πάνω στη γραμμή του.\par
\possindeo Το παλούκι είναι το κέντρο και το σχοινί η ακτίνα: κύκλος $(\mathrm{Ο},1{,}5)$ μ. Το δέντρο θα ήταν πάνω στον κύκλο αν $\mathrm{ΟΔ}=1{,}5$ μ.\par
\apantisi Ο κύκλος είναι ο $(\mathrm{Ο},1{,}5)$. Το δέντρο έχει $\mathrm{ΟΔ}=2\neq1{,}5$, άρα δεν είναι πάνω στη γραμμή.\par
\elegxos Το δέντρο απέχει περισσότερο από το σχοινί, άρα είναι έξω από το παρτέρι.

\tmima{Ασκήσεις}

\zoni{Το κατάλαβα}{σημεία, ευθείες, ημιευθείες, γωνίες}

\ex Συμπλήρωσε τα κενά. α) Το ίχνος της μύτης του μολυβιού μας δίνει την έννοια του \dots. β) Αν προεκτείνουμε ένα τμήμα και προς τα δύο άκρα του, προκύπτει \dots. γ) Από ένα σημείο διέρχονται \dots ευθείες, ενώ από δύο σημεία διέρχεται \dots ευθεία. δ) Αν προεκτείνουμε ένα τμήμα μόνο προς το ένα άκρο του, προκύπτει \dots. ε) Οι ημιευθείες που έχουν κοινή αρχή και σχηματίζουν ευθεία λέγονται \dots.

\ex Σχεδίασε τρία σημεία $\mathrm{Α},\,\mathrm{Β},\,\Gamma$ που δεν βρίσκονται στην ίδια ευθεία. Φέρε και ονόμασε τις ευθείες που ορίζουν και γράψε όλα τα ευθύγραμμα τμήματα που βλέπεις. Πόσες ευθείες και πόσα τμήματα είναι;

\ex Γράψε με σύμβολα: α) το τμήμα με άκρα $\mathrm{Κ}$ και $\Lambda$ \quad β) η ημιευθεία με αρχή το $\mathrm{Κ}$ που περνά από το $\Lambda$ \quad γ) η ευθεία που ορίζουν τα $\mathrm{Κ}$ και $\Lambda$.

\ex Σχεδίασε ένα σημείο $\mathrm{Ο}$ και δύο διαφορετικές ευθείες που να διέρχονται από αυτό. Ονόμασε τις αντικείμενες ημιευθείες με αρχή το $\mathrm{Ο}$. Πόσες ημιευθείες βλέπεις συνολικά;

\ex Σχεδίασε μια γωνία $x\hat{\mathrm{Ο}}y$ και γράψε το όνομά της με δύο διαφορετικούς τρόπους. Σημείωσε με χρώμα την κυρτή και με άλλο χρώμα τη μη κυρτή γωνία.

\ex Σχεδίασε με διαβήτη τον κύκλο $(\mathrm{Κ},2{,}5)$ cm και τρεις διαφορετικές ακτίνες του. Μέτρησέ τις. Τι παρατηρείς;

\ex Σωστό ή λάθος; α) Η ευθεία έχει δύο άκρα. \quad β) Τα τμήματα $\mathrm{ΑΒ}$ και $\mathrm{ΒΑ}$ είναι ίδια. \quad γ) Οι ημιευθείες $\mathrm{Αx}$ και $\mathrm{Βx}$ είναι ίδιες. \quad δ) Από δύο σημεία περνά μόνο μία ευθεία. \quad ε) Η γωνία έχει δύο ημιευθείες με κοινή αρχή. \quad στ) Οι ομόκεντροι κύκλοι έχουν την ίδια ακτίνα.

\zoni{Το χρησιμοποιώ}{ο κήπος, ο φακός, το στρογγυλό παρτέρι}

\ex \textbf{Η σειρά των δέντρων.} Ο παππούς Γιάννης θέλει να φυτέψει δέντρα σε ευθεία γραμμή. Πόσα παλούκια αρκεί να βάλει για να ορίσει την ευθεία και πώς θα την ελέγξει με ένα σχοινί;

\ex \textbf{Οι ακτίνες.} Ένας φακός φωτίζει με δέσμη που ξεκινά από το σημείο $\mathrm{Α}$ και περνά από το $\mathrm{Β}$. Ποια γεωμετρική έννοια περιγράφει η δέσμη; Αν δύο φακοί στέκονται αντικριστά και φωτίζουν ο ένας τον άλλον, τι σχέση έχουν οι δύο ημιευθείες;

\ex \textbf{Τα σχοινιά.} Πέντε παιδιά στέκονται σε σημεία του κήπου, από τα οποία τρία οποιαδήποτε δεν είναι στην ίδια ευθεία, και ενώνουν κάθε ζευγάρι με σχοινί. Πόσα σχοινιά χρειάζονται;

\ex \textbf{Το μονοπάτι.} Ένα ευθύγραμμο μονοπάτι $\varepsilon$ χωρίζει τον κήπο σε δύο ημιεπίπεδα. Τα δέντρα $\mathrm{Κ}$, $\Lambda$ είναι στην ίδια πλευρά και το $\mathrm{Μ}$ στην άλλη. Ποια από τα τμήματα $\mathrm{ΚΛ}$, $\mathrm{ΚΜ}$, $\Lambda\mathrm{Μ}$ τέμνουν το μονοπάτι;

\ex \textbf{Το παρτέρι.} Η Ρήνη χαράσσει με σχοινί $1{,}2$ μ. ένα στρογγυλό παρτέρι με κέντρο το παλούκι $\mathrm{Ο}$. Τρία φυτά απέχουν από το παλούκι $0{,}9$ μ., $1{,}2$ μ. και $1{,}6$ μ. Ποιο είναι πάνω στη γραμμή του κύκλου και ποιο εκτός του κύκλου;

\zoni{Το σκέφτομαι}{ευθεία ή τμήμα;}

\ex Γιατί η ευθεία δεν έχει μήκος αλλά το τμήμα έχει; Μπορείς να μετρήσεις το μήκος μιας ημιευθείας;

\ex Τρία σημεία $\mathrm{Α},\,\mathrm{Β},\,\Gamma$ βρίσκονται στην ίδια ευθεία με αυτή τη σειρά. Πόσες διαφορετικές ημιευθείες με αρχή ένα από αυτά τα σημεία ορίζονται πάνω στην ευθεία; (Μέτρα και τις δύο κατευθύνσεις.)

\ex Μπορούν δύο διαφορετικές ευθείες να περνούν και οι δύο από τα ίδια δύο σημεία $\mathrm{Α}$ και $\mathrm{Β}$; Εξήγησε με την ιδιότητα της ευθείας.

\ex Δύο ομόκεντροι κύκλοι έχουν ακτίνες $3$ cm και $5$ cm. Μπορεί να υπάρχει σημείο που ανήκει και στους δύο; Γιατί;

\zoni{Για το σπίτι}{γεωμετρία στο δωμάτιο}

\ex Βρες στο δωμάτιό σου πέντε παραδείγματα σημείου, ευθύγραμμου τμήματος και επιπέδου. Πού μοιάζει με ημιευθεία ένα αντικείμενο της καθημερινότητας;

\ex Με νήμα και δύο καρφίτσες σε ένα χαρτόνι δείξε ότι από δύο σημεία περνά μία ευθεία. Μετά σχεδίασε με διαβήτη έναν κύκλο που να περνά από ένα σημείο της επιλογής σου.
